\section{A first program}
\label{sec:basics:first_program}

A Ymir program is made of symbols. A symbol is a named entity declared in a
source file -- a function, a variable, a type -- and the name is what the rest of
the program uses to refer to it. Of the kinds of symbol, the one to know first is
the function, because every Ymir program contains at least one: the function
named \token{main}, which the system calls when the program starts.

The following listing is a complete program. It writes \textit{Hello World!} to
the console, and then stops.

\begin{lstlisting}[style=coloredverbatim, caption=Source file \textit{hello.yr}, label=lst:basics:hello_world]
use std::io;

fn main() {
  println("Hello World!");
}
\end{lstlisting}

The first line declares that the program uses something from the \token{std::io}
module of the standard library -- here the function \token{println}. Standard
library modules are available without any further installation step, so a
\token{use} line is all it takes to reach one.

The keyword \token{fn} introduces a function, and \token{main} names it. Its
body, between the braces, is a single statement: a call to \token{println}. The
parentheses after a function name hold the arguments passed to it, in this case
one string literal. Every statement ends with a semicolon. Once the body runs out
of statements, \token{main} returns, and the program ends with it.

\token{println} finishes the text it writes with a line break, so that whatever
is written next starts on a new line. The standard library also provides
\token{print}, which writes exactly what it is given and nothing more. A line
break can equally be written inside the literal itself, as the escape sequence
\tokennolst{\textbackslash{}n}: escape sequences start with a backslash and stand
for characters that cannot be typed directly in a literal.
Chapter~\ref{chap:scalars} lists them. Note that a string literal is enclosed in
double quotes; single quotes are reserved for character literals, described in
that same chapter.

\subsection{Compilation and execution}
\label{sec:basics:compilation_execution}

A source file is not yet a program. The compiler, \token{gyc}, turns it into one:
given \textit{hello.yr}, the command below produces an executable file named
\textit{hello}.

\begin{lstlisting}[style=bashVerb, escapechar=@+]
@+\textcolor{teal!80}{alice@dev:\mtilde}@+$ gyc hello.yr -o hello
\end{lstlisting}

That file now sits in the current directory, and running it prints the expected
line.

\begin{lstlisting}[style=bashVerb, escapechar=@+]
@+\textcolor{teal!80}{alice@dev:\mtilde}@+$ ./hello
Hello World!
\end{lstlisting}


\vfill
\pagebreak
